\documentclass[11pt,reqno]{amsart}
\usepackage[localtheorems]{raeez-math-template}

% ================================================================
% Packages
% ================================================================
\usepackage{array}

% ================================================================
% Theorem environments
% ================================================================
\newtheorem{theorem}{Theorem}[section]
\newtheorem{proposition}[theorem]{Proposition}
\newtheorem{lemma}[theorem]{Lemma}
\newtheorem{corollary}[theorem]{Corollary}
\newtheorem{conjecture}[theorem]{Conjecture}
\theoremstyle{definition}
\newtheorem{definition}[theorem]{Definition}
\newtheorem{construction}[theorem]{Construction}
\newtheorem{example}[theorem]{Example}
\theoremstyle{remark}
\newtheorem{remark}[theorem]{Remark}

% ================================================================
% Macros
% ================================================================
\providecommand{\cA}{\mathcal{A}}
\providecommand{\cD}{\mathcal{D}}
\providecommand{\cH}{\mathcal{H}}
\providecommand{\cM}{\mathcal{M}}
\providecommand{\cW}{\mathcal{W}}
\providecommand{\barB}{\overline{B}}
\providecommand{\Ran}{\mathrm{Ran}}
\providecommand{\FM}{\overline{\mathrm{FM}}}
\providecommand{\Conf}{\mathrm{Conf}}
\providecommand{\fg}{\mathfrak{g}}
\providecommand{\Sym}{\mathrm{Sym}}
\providecommand{\Res}{\operatorname{Res}}
\providecommand{\Spec}{\operatorname{Spec}}
\providecommand{\Diff}{\mathrm{Diff}}
\providecommand{\Vir}{\mathrm{Vir}}
\providecommand{\Hom}{\mathrm{Hom}}
\providecommand{\id}{\mathrm{id}}
\providecommand{\ChirHoch}{\mathrm{ChirHoch}}
\renewcommand{\Bbbk}{\mathbf{C}}

\DeclareMathOperator{\gr}{gr}
\DeclareMathOperator{\End}{End}

\numberwithin{equation}{section}

\makeatletter
\let\ps@oldref\ref
\renewcommand{\ref}[1]{%
 \@ifundefined{r@#1}{\textit{full monograph}}{\ps@oldref{#1}}%
}
\makeatother

% ================================================================
\begin{document}

\title[Programme Summary: Sections 2--4]
{Modular Koszul Duality: Sections 2--4}

\author{Raeez Lorgat}
\address{Perimeter Institute for Theoretical Physics,
 Waterloo, ON N2L 2Y5, Canada}
\email{rlorgat@perimeterinstitute.ca}

\date{\today}

\begin{abstract}
Sections~2--4 develop the curvature calculus of the modular bar
construction. Section~2 installs the genus-$g$ propagator via the
prime form, computes the period defect, and identifies the single
scalar $\kappa(\cA)$ controlling fibrewise curvature
$d_{\mathrm{fib}}^2 = \kappa\cdot\omega_g$. Section~3 develops the
modular homotopy theory: the convolution $L_\infty$-algebra
$(\mathfrak{g}^{\mathrm{mod}}_\cA, \mu_r)$, the modular operad with
Getzler--Kapranov Feynman transform, the six-component total
differential, and the logarithmic Fulton--MacPherson compactification
$\overline{\FM}^{\log}$ of Mok~25 as ambient. Section~4 constructs
the universal Maurer--Cartan element
$\Theta_\cA = D_\cA - d_0 \in \mathfrak{g}^{\mathrm{mod}}_\cA$,
proves $[\Theta_\cA, \Theta_\cA] = 0$ from $D_\cA^2 = 0$, and
records the tridegree, depth, and convergence discipline. The five
theorems are projections of this single universal object.
\end{abstract}

\maketitle

% ====================================================================
\section{Curvature and the genus tower}
\label{sec:curvature}
% ====================================================================

At genus~$0$ the bar complex is flat: $d^2 = 0$ is an
algebraic identity, equivalent to the Borcherds axiom.
On a family of curves of genus $g \ge 1$, this nilpotence
fails. The failure is controlled by a single scalar.

% ====================================================================
\subsection{The prime form and the period defect}
\label{ssec:prime-form}
% ====================================================================

On a compact Riemann surface~$\Sigma_g$ of genus~$g \ge 1$,
the function $z_1 - z_2$ has no global meaning: a
meromorphic function with a single simple zero must have a
simple pole. The \emph{prime form} $E(z_1, z_2)$, a
section of $K^{-1/2} \boxtimes K^{-1/2}$ vanishing to
first order along the diagonal, replaces it. The
genus-$g$ propagator
\begin{equation}\label{eq:genus-g-propagator}
\eta^{(g)}_{12}
\;=\;
d\log E(z_1, z_2)
+ \pi \sum_{\alpha,\beta = 1}^{g}
(\operatorname{Im}\Omega)^{-1}_{\alpha\beta}
\,\omega_\alpha(z_1)\,\bar\omega_\beta(z_2)
\end{equation}
is single-valued; the period correction couples to the
Hodge bundle $\mathbb{E} = R^0\pi_*\omega_\pi$ through
the period matrix~$\Omega$.

Three properties of~\eqref{eq:genus-g-propagator} constrain
the genus tower. First, the logarithmic derivative
$d\log E$ has conformal weight~$1$ in both variables,
regardless of the conformal weight of the field being
sewed: every edge of every graph sum at every genus uses
the standard Hodge bundle~$\mathbb{E}_1$. Second, the
period correction is harmonic in each variable and involves
$(\operatorname{Im}\Omega)^{-1}$, the Arakelov Green
function on~$\Sigma_g$. Third, the Arnold relation acquires
a defect proportional to the Arakelov form, and this defect
propagates into the bar differential.

% ====================================================================
\subsection{Curvature}
\label{ssec:curvature}
% ====================================================================

Let $\pi \colon \mathcal{C}_g \to \overline{\cM}_g$ be the
universal curve. On each fibre~$\Sigma_g$, the propagator
$\eta^{(g)}$ is single-valued but no longer closed: the
period correction has a $\bar\partial$-component, and the
Arnold relation fails. The fibrewise bar differential
satisfies
\begin{equation}\label{eq:d-squared-curvature}
d_{\mathrm{fib}}^{\,2}
\;=\;
\kappa(\cA) \cdot \omega_g\,,
\end{equation}
where $\omega_g$ is the Arakelov $(1,1)$-form on the fibre
and $\kappa(\cA) \in \Bbbk$ is the \emph{modular
characteristic}. The modular characteristic is intrinsic to
the chiral algebra; it is determined by the leading OPE
singularity and is visible already at genus~$0$ through the chain
\[
K_\cA^{\mathrm{coll}}\longmapsto
q_2(K_\cA^{\mathrm{coll}})\longmapsto
\tau_{\cA,2}^{\mathrm{res}}q_2(K_\cA^{\mathrm{coll}}).
\]
The last term is $\kappa(\cA)$ for abelian and Virasoro-type
families. For non-abelian affine Kac--Moody in the trace-form
convention it is $\kappa_{\mathrm{dp}}(\cA)$ and
$\kappa(\cA) = \kappa_{\mathrm{dp}}(\cA) + \dim(\fg)/2$.

The failure of $d^2 = 0$ at genus~$g \ge 1$ is not a defect
of the construction; it is the construction detecting global
geometry. Period integrals restore nilpotence: the
Gauss--Manin connection absorbs the fibrewise curvature, and
the total differential satisfies $D_g^2 = 0$. The
curvature~\eqref{eq:d-squared-curvature} is the
infinitesimal avatar of this restoration.

% ====================================================================
\subsection{The genus expansion}
\label{ssec:genus-expansion}
% ====================================================================

For a modular Koszul chiral algebra~$\cA$ and $g\geq2$, Theorem~D
begins with
$\operatorname{Obs}^{\mathrm{def}}_g(\cA)
\in H^2(\operatorname{Def}_g(\cA))$.  Under $H_D^K$,
\begin{equation}\label{eq:genus-universality}
\mathfrak O_g^K(\cA)=\kappa(\cA)\lambda_{-1}(\mathbb E_g),
\qquad
\operatorname{ch}_g(\mathfrak O_g^K(\cA))
=(-1)^g\kappa(\cA)\lambda_g.
\end{equation}
On the scalar-diagonal graph lane the free energy is
$F_g^{\mathrm{sc}}(\cA)=\kappa(\cA)\lambda_g^{\mathrm{FP}}$, where
\begin{equation}\label{eq:faber-pandharipande}
\lambda_g^{\mathrm{FP}}
\;=\;
\frac{2^{2g-1} - 1}{2^{2g-1}}
\cdot
\frac{|B_{2g}|}{(2g)!}
\end{equation}
is the Faber--Pandharipande coefficient ($B_{2g}$ the
$2g$-th Bernoulli number), packages the entire genus tower
into the Hirzebruch $\hat{A}$-class:
\begin{equation}\label{eq:ahat-generating}
\sum_{g=1}^{\infty}
F_g^{\mathrm{sc}}(\cA)\,\hbar^{2g-2}
\;=\;
\frac{\kappa(\cA)}{\hbar^2}
\bigl[\hat{A}(i\hbar) - 1\bigr]\,,
\qquad
\hat{A}(x) = \frac{x/2}{\sinh(x/2)}\,.
\end{equation}
At genus~$1$: $\lambda_1^{\mathrm{FP}} = 1/24$. At
genus~$2$: $\lambda_2^{\mathrm{FP}} = 7/5760$. The
universal ratio $F_2/F_1 = 7/240$ is independent
of~$\cA$. The series converges for
$|\hbar| < 2\pi$, the first zero of $\sin(\hbar/2)$.

At genus~$1$, the pointed trace identity
$\operatorname{tr}_1\operatorname{Obs}^{\mathrm{def}}_{1,1}
=\kappa\lambda_1$ holds for every modular Koszul family.  At
$g\ge2$, the graph output carries the cross-channel decomposition
of Remark~\ref{rem:multi-weight}.

\begin{remark}[Multi-weight correction]
\label{rem:multi-weight}
For a chiral algebra with generators of distinct conformal
weights $h_1, \ldots, h_r$, the genus-$g$ free energy
decomposes as
\begin{equation}\label{eq:multi-weight-decomposition}
F_g(\cA)
\;=\;
\kappa(\cA) \cdot \lambda_g^{\mathrm{FP}}
\;+\;
\delta F_g^{\mathrm{cross}}(\cA)\,,
\end{equation}
where $\delta F_g^{\mathrm{cross}}$ is a graph sum over
mixed-channel boundary graphs of~$\overline{\cM}_{g,0}$.
At genus~$1$, $\delta F_1^{\mathrm{cross}} = 0$
unconditionally. For uniform-weight algebras
($h_1 = \cdots = h_r$),
$\delta F_g^{\mathrm{cross}} = 0$ at all genera. For
multi-weight algebras at $g \ge 2$, the correction is
generically nonzero: for $\mathcal{W}_3$ at genus~$2$,
$\delta F_2^{\mathrm{cross}}(\mathcal{W}_3)
= (c + 204)/(16c) > 0$ for all $c > 0$.
\end{remark}

% ====================================================================
\subsection{The modular characteristic}
\label{ssec:kappa}
% ====================================================================

The modular characteristic $\kappa(\cA)$ is the single
scalar controlling the genus tower. It is computed from
the leading OPE singularity and fixed universally by the
genus-$1$ curvature extracted through the bar construction.

\begin{table}[ht]
\centering
\caption{Modular characteristic for the standard families}
\label{tab:kappa}
\renewcommand{\arraystretch}{1.6}
\begin{tabular}{lcc}
\toprule
\textbf{Algebra $\cA$}
 & $\boldsymbol{\kappa(\cA)}$
 & \textbf{Shadow depth class} \\
\midrule
Heisenberg $\cH_k$
 & $k$
 & $\mathbf{G}$ \\
Affine Kac--Moody $\widehat{\fg}_k$
 & $\dfrac{\dim(\fg) \cdot (k + h^\vee)}{2h^\vee}$
 & $\mathbf{L}$ \\[4pt]
$\beta\gamma$ system
 & $c/2$
 & $\mathbf{C}$ \\
Virasoro $\Vir_c$
 & $c/2$
 & $\mathbf{M}$ \\
$\mathcal{W}_3^k(\mathfrak{sl}_3)$
 & $5c/6$
 & $\mathbf{M}$ \\
$\mathcal{W}_N^k(\mathfrak{sl}_N)$
 & $\displaystyle c \cdot \sum_{j=2}^{N}\frac{1}{j}$
 & $\mathbf{M}$ \\[4pt]
\bottomrule
\end{tabular}
\end{table}

The modular characteristic is additive under independent
sums ($\kappa(\cA \oplus \cA') = \kappa(\cA) + \kappa(\cA')$
when the mixed OPE vanishes) and constrained by Koszul
duality: $\kappa(\cA) + \kappa(\cA^!) = 0$ for
Kac--Moody and free-field algebras;
$\kappa(\Vir_c) + \kappa(\Vir_{26-c}) = 13$ for Virasoro.
In general, $\kappa(\cA) + \kappa(\cA^!) = \rho \cdot K$
for $\mathcal{W}$-algebras, where $\rho$ is the
anomaly ratio and $K$ the complementarity constant.

At the self-dual point ($c = 13$ for Virasoro), the
complementarity sum is symmetric:
$\kappa(\Vir_{13}) = 13/2 = \kappa(\Vir_{13}^!)$.
At the critical level ($k = -h^\vee$ for Kac--Moody),
$\kappa = 0$ and the bar complex is flat at all genera:
the content migrates from curvature to bar cohomology,
where it appears as the Feigin--Frenkel centre.


% ====================================================================
\section{The five main theorems}
\label{sec:five-theorems}
% ====================================================================

The bar construction of a chiral algebra on a curve produces
a factorization coalgebra on the Ran space of the curve.
The five main theorems describe what this coalgebra
determines: the Koszul dual, the original algebra,
the complementarity decomposition, the genus tower, and
the chiral Hochschild cohomology. Each theorem is a
projection of the universal Maurer--Cartan element
$\Theta_\cA$ (Section~\ref{sec:mc}).

% ====================================================================
\subsection{Theorem A: Verdier intertwining}
\label{ssec:thm-a}
% ====================================================================

\begin{theorem}[A]\label{thm:A}
Let\/ $\cA$ be a chiral algebra on a smooth projective
curve~$X$. There is an adjunction
\[
B
\;\colon\;
\mathrm{Alg}_{\mathrm{aug}}(\mathrm{Fact}(X))
\;\rightleftarrows\;
\mathrm{CoAlg}_{\mathrm{conil}}(\mathrm{Fact}(X))
\;\colon\;
\Omega
\]
between augmented factorization algebras and conilpotent
factorization coalgebras, with unit and counit the
universal twisting morphisms.  On the dualizable Ran surface,
Verdier duality sends the bar coalgebra to the factorization algebra
\begin{equation}\label{eq:verdier}
K_X(\cA):=\mathbb{D}_{\Ran}B(\cA).
\end{equation}
For a quadratic presentation, the comparison
$q_\cA\colon\cA^{\mathrm i}\to B(\cA)$ has Verdier transpose
\[
\mathbb D(q_\cA)\colon K_X(\cA)\longrightarrow
\mathbb D_{\Ran}(\cA^{\mathrm i}).
\]
A chosen Koszul pair supplies an algebra map
$\nu_\cA\colon K_X(\cA)\to\cA^!$, an equivalence under the
finite-type or completed pair hypotheses.
\end{theorem}

The adjunction uses the augmented factorization-algebra structure.
The Verdier clause additionally uses dualizability and monoidal
exchange on~$\Ran(X)$.  Equation~\eqref{eq:verdier} retains the full
chain-level data; the pair map compares it with the strict algebra
$\cA^!=(H^*(B(\cA)))^\vee$.

% ====================================================================
\subsection{Theorem B: bar-cobar inversion}
\label{ssec:thm-b}
% ====================================================================

\begin{theorem}[B]\label{thm:B}
On the Koszul locus, the cobar of the bar recovers the
original algebra:
\begin{equation}\label{eq:inversion}
\Omega\bigl(B(\cA)\bigr)
\;\xrightarrow{\;\sim\;}
\;\cA\,.
\end{equation}
The counit of the bar-cobar adjunction is a
quasi-isomorphism of augmented factorization algebras.
At genus~$0$ this is unconditional. For arbitrary
modular pre-Koszul inputs at $g \ge 1$, the ordinary direct-sum
statement holds only on the Koszul/collapse loci. Off that locus,
and for class~$M$, inversion is a weight-completed coderived
Positselski statement.
\end{theorem}

Bar-cobar inversion is a round-trip, not a duality.
It does not produce the Koszul dual (that is Verdier
duality, Theorem~A); it does not produce the bulk
algebra (that is the chiral derived centre, Theorem~H).
It recovers~$\cA$ itself.

The proof proceeds by constructing a contracting homotopy
for the counit, using the PBW filtration as a control
parameter: the associated graded of the cobar complex
coincides with a cofree resolution of the PBW-graded
algebra, and the filtration spectral sequence degenerates
at~$E_2$. Convergence of the spectral sequence requires
the Koszul hypothesis; for non-Koszul algebras, bar-cobar
inversion requires completion.

% ====================================================================
\subsection{Theorem C: complementarity}
\label{ssec:thm-c}
% ====================================================================

\begin{theorem}[C]\label{thm:C}
For a Koszul pair $(\cA, \cA^!)$ on~$X$, Theorem~C has three typed
tiers.  \textsc{C0} is the strict flat fibre-centre comparison giving
the ambient complex
\[
\mathbf{C}_g(\cA)
\;:=\;
R\Gamma\bigl(\overline{\cM}_g,\, \mathcal{Z}_\cA\bigr)
\]
on the chirally Koszul datum.  \textsc{C1}, once the represented
Verdier involution, perfectness, and nondegenerate anti-invariant
Verdier pairing are supplied, decomposes this complex into homotopy
eigenspaces:
\begin{equation}\label{eq:complementarity}
\mathbf{C}_g
\;\simeq\;
\mathbf{Q}_g(\cA) \;\oplus\; \mathbf{Q}_g(\cA^!)\,.
\end{equation}
For $g \ge 1$ the pairing has cohomological degree $-(3g-3)$, and
the summands $\mathbf{Q}_g(\cA)=
\operatorname{fib}(\sigma-\mathrm{id})$ and
$\mathbf{Q}_g(\cA^!)=\operatorname{fib}(\sigma + \mathrm{id})$
are Lagrangian under the C1 hypotheses.  \textsc{C2} is the
shifted-symplectic/BV upgrade; it is conditional on the C2 hypothesis
package and is not supplied by the scalar conductor or by
uniform-weight diagonalization alone.
\end{theorem}

The scalar identity $F_g=\kappa\lambda_g^{\mathrm{FP}}$ is the
uniform-weight scalar trace shadow of C1 together with Theorem~D.  It
is not the C2 shifted-symplectic upgrade.  The C2 statement requires
the Pantev--To\"en--Vaqui\'e--Vezzosi shifted-symplectic package,
cyclic closedness, skew-adjointness, modular-bracket compatibility,
twisted-tensor acyclicity, and the stated BV/QME anomaly-killing
inputs.

% ====================================================================
\subsection{Theorem D: the four-stage modular trace}
\label{ssec:thm-d}
% ====================================================================

\begin{theorem}[D]\label{thm:D}
Let~$\cA$ be a cyclic chart algebra with modular characteristic
$\kappa(\cA)$.  For $g\geq2$, its native genus-$g$ class is
$\operatorname{Obs}^{\mathrm{def}}_g(\cA)\in
H^2(\operatorname{Def}_g(\cA))$.  Under~$H_D^1$,
$\operatorname{tr}_1\operatorname{Obs}^{\mathrm{def}}_{1,1}(\cA)
=\kappa(\cA)\lambda_1$.  Under~$H_D^K$,
\[
 \mathfrak O_g^K(\cA)=\kappa(\cA)\lambda_{-1}(\mathbb E_g),
 \qquad
 \operatorname{ch}_g(\mathfrak O_g^K(\cA))
 =(-1)^g\kappa(\cA)\lambda_g.
\]
The package~$H_D^{\mathrm{tr}}$ supplies the deformation-to-$K$
comparison.  Under~$H_D^{\mathrm{graph}}$,
\[
 F_g(\cA)=\kappa(\cA)\lambda_g^{\mathrm{FP}}
 +\delta F_g^{\mathrm{cross}}(\cA).
\]
The scalar-diagonal graph package gives
$\delta F_g^{\mathrm{cross}}(\cA)=0$.
\end{theorem}

The $\hat A$-series evaluates the scalar-diagonal graph functional.
The native class, the perfect object, the Hodge character, and this
number lie in different ambient groups; $H_D^{\mathrm{tr}}$ is the
map between the first two stages.

% ====================================================================
\subsection{Theorem H: chiral Hochschild cohomology}
\label{ssec:thm-h}
% ====================================================================

\begin{theorem}[H]\label{thm:H}
Fix a finite set $S\subset\mathbb Z$.  A family datum
$H_H(\cA;S)$ consists of a complete chart complex $Q_\cA$, a
quasi-isomorphism to $C^\bullet_{\mathrm{ch}}(\cA,\cA)$, and a
strong deformation retract
\[
 K_{\cA,S}\mathop{\xrightarrow{\ \iota\ }}Q_\cA
 \mathop{\xrightarrow{\ p\ }}K_{\cA,S},
 \qquad p\iota=\id,
 \qquad dh+hd=\id-\iota p,
\]
including the finite-window completion and ordered-to-symmetric
comparison.  Then
\[
 \ChirHoch^n(\cA)\cong H^n(K_{\cA,S}),
 \qquad
 \operatorname{Supp}\ChirHoch^\bullet(\cA)\subseteq S.
\]
For finite-dimensional graded pieces, the Hilbert polynomial is
$P_\cA(t)=\sum_{n\in S}\dim H^n(K_{\cA,S})t^n$.  A perfect
degree-$d$ pairing with $\cA^!$ is an additional datum.
\end{theorem}

The chiral derived centre
$Z^{\mathrm{der}}_{\mathrm{ch}}(\cA) =
C^{\bullet}_{\mathrm{ch}}(\cA, \cA)$ is the universal
closed-sector algebra: the cochain-level operator algebra acting on
the boundary algebra~$\cA$. A physical bulk interpretation requires
the open--closed comparison datum. Theorem~H computes its cohomology
from $K_{\cA,S}$.  Bakalov--De Sole--Kac compute the bounded
Virasoro groups as one-dimensional in degrees $\{0,2,3\}$
\cite[Theorem~7.2]{BDSK21}; a
bounded-to-chart quasi-isomorphism transports that vector to the
curve chart.

% ====================================================================
\subsection{The Koszulness meta-theorem}
\label{ssec:meta-theorem}
% ====================================================================

The hypothesis shared by Theorems~A--H is chiral
Koszulness: bar cohomology $H^*(B(\cA))$ concentrated in
bar degree~$1$. The meta-theorem organises this condition by
twelve tests and consequences: seven independent bidirectional
equivalences, one listed consequence, two conditional
comparisons, one one-way chiral Hochschild consequence, and one
one-directional purity implication.

\begin{theorem}[Koszulness characterization]
\label{thm:koszul-meta}
For a chiral algebra~$\cA$ on~$X$ with PBW filtration,
conditions \textup{(i)--(v)} and \textup{(ix)--(x)} are the
seven independent bidirectional equivalences. Item
\textup{(vi)} follows from \textup{(v)}. Item \textup{(vii)}
is conditional on the factorization-homology comparison and
detection package. Item \textup{(viii)} is a one-way consequence
of a family datum $H_H(\cA;S)$.
Under the shifted-symplectic hypothesis package, \textup{(xi)} is
equivalent. Condition \textup{(xii)} implies \textup{(x)}; the
converse is open.
\begin{enumerate}[label=\textup{(\roman*)}]
\item $\cA$ is chirally Koszul.
\item The PBW spectral sequence collapses at\/~$E_2$.
\item The minimal $A_\infty$-model of $B(\cA)$
 is formal: $m_n = 0$ for $n \ge 3$.
\item Ext diagonal vanishing:
 $\operatorname{Ext}^{p,q}_\cA(\omega_X, \omega_X) = 0$
 for $p \ne q$.
\item The bar-cobar counit
 $\Omega(B(\cA)) \to \cA$ is a quasi-isomorphism.
\item The Barr--Beck--Lurie comparison is an equivalence.
\item Factorization homology
 $\int_{\Sigma_g}\!\cA$ is concentrated in degree~$0$
 for all~$g$.
\item Under $H_H(\cA;S)$,
 $\operatorname{Supp}\mathrm{ChirHoch}^\bullet(\cA)\subseteq S$,
 with Hilbert polynomial
 $\sum_{n\in S}\dim H^n(K_{\cA,S})t^n$ when the graded pieces
 are finite dimensional.
\item The Kac--Shapovalov determinant is nonzero in the
 bar-relevant range.
\item Restriction to every FM boundary stratum is
 acyclic outside degree~$0$.
\item[{\upshape (xi)}] Lagrangian criterion \textup{(conditional)}.
\item[{\upshape (xii)}] $\cD$-module purity
 \textup{(one-directional: (x) $\Rightarrow$ (xii))}.
\end{enumerate}
\end{theorem}

The core circuit of independent equivalences runs through PBW
degeneration (commutative algebra), $A_\infty$-formality
(homotopical algebra), Ext diagonal vanishing, bar-cobar
inversion, Kac--Shapovalov nonvanishing (representation theory),
and FM boundary acyclicity (algebraic topology). The
Barr--Beck--Lurie row records the higher-categorical consequence
of bar-cobar inversion; the factorization-homology row needs its
comparison and detection hypotheses.


% ====================================================================
\section{The universal Maurer--Cartan element
 and the shadow tower}
\label{sec:mc}
% ====================================================================

The five theorems of Section~\ref{sec:five-theorems} are
projections of a single algebraic object: the universal
Maurer--Cartan element $\Theta_\cA$ of the modular
convolution algebra.

% ====================================================================
\subsection{The bar-intrinsic construction}
\label{ssec:bar-intrinsic}
% ====================================================================

The genus-completed bar differential
$D_\cA = \sum_{g \ge 0} \hbar^g\,d^{(g)}_\cA$
satisfies $D_\cA^2 = 0$ by the codimension-$2$ face
cancellation on $\overline{\cM}_{g,n}$ (the same mechanism
as~\S\ref{ssec:d-squared}, extended to all genera by Mok's
log Fulton--MacPherson compactification). Write
$d_0 = d^{(0)}_\cA$ for the genus-$0$ bar differential.

\begin{construction}[Bar-intrinsic MC element]
\label{constr:bar-intrinsic}
Define the positive-genus correction
\begin{equation}\label{eq:theta-def}
\Theta_\cA
\;:=\;
D_\cA - d_0
\;=\;
\sum_{g \ge 1}
\hbar^g\,d^{(g)}_\cA
\;\in\;
\prod_{g \ge 1}
\mathrm{CoDer}^{\mathrm{cyc}}
\bigl(B^{(g)}_X(\cA)\bigr)[1]\,.
\end{equation}
\end{construction}

\begin{theorem}\label{thm:theta-mc}
$\Theta_\cA$ is a Maurer--Cartan element of the modular
cyclic deformation complex:
\begin{equation}\label{eq:mc-equation}
d_0(\Theta_\cA)
+ \tfrac{1}{2}[\Theta_\cA, \Theta_\cA]
\;=\; 0\,.
\end{equation}
\end{theorem}

\begin{proof}[Proof sketch]
Expand $D_\cA^2 = (d_0 + \Theta_\cA)^2 = 0$.
Since $d_0^2 = 0$ (genus-$0$ bar flatness), the
remaining terms give
$d_0(\Theta_\cA) + \frac{1}{2}[\Theta_\cA, \Theta_\cA]
= 0$, which is the MC equation. The bracket is the
convolution bracket on the modular cyclic deformation
complex, inherited from stable-curve gluing; its square-zero
property is a formal consequence of $d_0^2 = 0$.
\end{proof}

The construction requires no restriction to simple Lie
symmetry, no one-channel hypothesis, and no tautological-line
support condition. It works for all modular Koszul chiral
algebras. The modular characteristic~$\kappa$, the cubic
shadow~$\mathfrak{C}$, and the quartic resonance
class~$\mathfrak{Q}$ are degree projections of this single
element:
\begin{equation}\label{eq:shadow-projections}
\kappa(\cA) = \pi_2(\Theta_\cA)\,,
\qquad
\mathfrak{C}(\cA) = \pi_3(\Theta_\cA)\,,
\qquad
\mathfrak{Q}(\cA) = \pi_4(\Theta_\cA)\,.
\end{equation}
Each degree projection carries more information than the
last; the full $\Theta_\cA$ carries all of it.

% ====================================================================
\subsection{The shadow obstruction tower}
\label{ssec:shadow-tower}
% ====================================================================

The degree-$r$ projection
$\mathrm{Sh}_r(\cA) := \pi_r(\Theta_\cA)$ is the $r$-th
shadow coefficient. The sequence $S_2 = \kappa$,
$S_3 = \alpha$, $S_4$, $S_5, \ldots$ constitutes the
\emph{shadow obstruction tower}. On any one-dimensional
primary slice~$L$ of the cyclic deformation complex, the
MC equation~\eqref{eq:mc-equation} imposes a quadratic
constraint.

\begin{theorem}[Riccati algebraicity]
\label{thm:riccati}
The weighted shadow generating function
$H(t) := \sum_{r \ge 2} r\,S_r\,t^r$
satisfies
\begin{equation}\label{eq:algebraic-relation}
H(t)^2 \;=\; t^4\,Q_L(t)\,,
\end{equation}
where
\begin{equation}\label{eq:shadow-metric}
Q_L(t)
\;=\;
(2\kappa + 3\alpha t)^2
\;+\;
2\Delta\,t^2\,,
\qquad
\Delta := 8\kappa S_4\,,
\end{equation}
is the \emph{shadow metric}: a quadratic polynomial in~$t$
determined by three genus-$0$ invariants
$(\kappa, \alpha, S_4)$.
The entire shadow obstruction tower on~$L$ is determined
by~$Q_L$.
\end{theorem}

The generating function $H(t) = t^2\sqrt{Q_L(t)}$ is
algebraic of degree~$2$. The shadow obstruction tower
is neither a Taylor expansion that might diverge nor a
formal power series requiring analytic continuation: it is
the expansion of an explicit algebraic function. The
MC equation, which a priori is an infinite system of coupled
nonlinear equations, reduces on each primary slice to a
single quadratic identity.

% ====================================================================
\subsection{The G/L/C/M classification}
\label{ssec:glcm}
% ====================================================================

The discriminant $\Delta = 8\kappa S_4$ of the shadow
metric forces a universal dichotomy.

\begin{definition}[Shadow depth classification]
\label{def:shadow-depth}
The \emph{shadow depth} of~$\cA$ is
$r_{\max}(\cA) := \sup\{r \ge 2 : S_r \ne 0\}$.
The shadow depth class is:
\begin{center}
\renewcommand{\arraystretch}{1.2}
\begin{tabular}{clcl}
\toprule
\textbf{Class} & \textbf{Name} &
 $\boldsymbol{r_{\max}}$ & \textbf{Mechanism} \\
\midrule
$\mathbf{G}$ & Gaussian &
 $2$ & $Q_L$ perfect square, $\alpha = 0$ \\
$\mathbf{L}$ & Lie &
 $3$ & $Q_L$ perfect square, $\alpha \ne 0$ \\
$\mathbf{C}$ & Contact &
 $4$ & stratum separation \\
$\mathbf{M}$ & Mixed &
 $\infty$ & $\Delta \ne 0$: $Q_L$ irreducible \\
\bottomrule
\end{tabular}
\end{center}
\end{definition}

The mechanism: when $\Delta = 0$, the shadow metric is a
perfect square $Q_L(t) = (2\kappa + 3\alpha t)^2$ and
$\sqrt{Q_L}$ is polynomial: the tower terminates.
If $\alpha = 0$ it terminates at degree~$2$ (class~$\mathbf{G}$:
Heisenberg); if $\alpha \ne 0$ at degree~$3$
(class~$\mathbf{L}$: affine Kac--Moody). When
$\Delta \ne 0$, $\sqrt{Q_L}$ is irrational: the tower
runs forever (class~$\mathbf{M}$: Virasoro,
$\mathcal{W}_N$). Class~$\mathbf{C}$ ($\beta\gamma$)
escapes the dichotomy by stratum separation: the quartic
contact invariant lives on a charged stratum whose
self-bracket exits the complex by rank-one rigidity.

The four classes are a homotopy-invariant classification:
they coincide with the $L_\infty$-formality depth of the
transferred minimal model. Class~$\mathbf{G}$ is
$L_\infty$-formal (all higher brackets vanish).
Class~$\mathbf{L}$ has exactly one nontrivial Massey
product. Class~$\mathbf{C}$ has formality obstructions
through degree~$4$. Class~$\mathbf{M}$ is intrinsically
non-formal: no finite truncation of the $L_\infty$
structure is quasi-isomorphic to a dg~Lie algebra.

Archetypes:
\begin{center}
\renewcommand{\arraystretch}{1.3}
\begin{tabular}{clll}
\toprule
\textbf{Class} & \textbf{Archetype}
 & $\boldsymbol{\kappa}$ & $\boldsymbol{\Delta}$ \\
\midrule
$\mathbf{G}$ & Heisenberg $\cH_k$
 & $k$ & $0$ \\
$\mathbf{L}$ & $\widehat{\fg}_k$
 & $\dim(\fg)(k{+}h^\vee)/(2h^\vee)$ & $0$ \\
$\mathbf{C}$ & $\beta\gamma$
 & $c/2$ & $0$ (stratum separation) \\
$\mathbf{M}$ & $\Vir_c$
 & $c/2$ & $40/(5c{+}22)$ \\
\bottomrule
\end{tabular}
\end{center}

For the Virasoro algebra, the quartic contact invariant is
$Q^{\mathrm{contact}}_{\Vir} = 10/[c(5c+22)]$ and the
critical discriminant is $\Delta_{\Vir} = 8 \cdot (c/2)
\cdot 10/[c(5c+22)] = 40/(5c+22)$. Every shadow
coefficient $S_r$ for $r \ge 5$ is a rational function
of~$c$ with poles only at $c = 0$ and $c = -22/5$:
\begin{align}
S_5 &= -\frac{48}{c^2(5c{+}22)}\,,
 &
S_6 &= \frac{80(45c{+}193)}{3\,c^3(5c{+}22)^2}\,,
 &
S_7 &= -\frac{2880(15c{+}61)}{7\,c^4(5c{+}22)^2}\,.
\label{eq:shadow-coefficients}
\end{align}
These are the fingerprints of a class-$\mathbf{M}$ algebra:
the shadow tower is infinite, alternating in sign, and
controlled by an irrational algebraic function.

% ====================================================================
\subsection{The gauge-gravity dichotomy}
\label{ssec:gauge-gravity}
% ====================================================================

The classification has a physical interpretation. At
$\kappa = 0$ the bar complex is flat at all genera:
the genus tower carries no curvature, the free energy
vanishes, and the content lives in bar cohomology. This
is the regime of gauge theory: the Feigin--Frenkel centre
(the Langlands dual of the affine algebra at critical level)
is the genus-$0$ bar cohomology $H^0(B(\widehat{\fg}_{-h^\vee}))$.

At $\kappa \ne 0$ the bar complex is curved: every genus
contributes a nonzero obstruction class, the genus tower
is nontrivial, and the free energy is the $\hat{A}$-genus.
This is the regime of gravity: the bar curvature is the
perturbative genus expansion.

The Koszul duality $\cA \leftrightarrow \cA^!$ exchanges
gauge and gravity sectors. For affine Kac--Moody at
generic level, $\kappa(\widehat{\fg}_k) \ne 0$ and the
algebra is curved; at the Koszul-dual critical level
$k' = -k - 2h^\vee$, $\kappa(\widehat{\fg}_{k'}) \ne 0$
as well, and both sides are gravitational. The gauge regime
(critical level, $\kappa = 0$) is the fixed locus of a
one-parameter family.

For the Virasoro algebra, the Koszul dual is
$\Vir_c^! = \Vir_{26-c}$. The gauge locus is $c = 0$
(topological: $\kappa = 0$); the gravitational locus is
$c \ne 0$. The self-dual point $c = 13$
($\kappa = 13/2 = \kappa^!$) is the critical string.
The Clifford algebra of the complementarity pairing
degenerates at $c = 26$ ($\kappa(\Vir_0^!) = 0$):
the $26$-dimensional bosonic string is the critical point
where the Koszul dual becomes flat.

The dichotomy gauge/$\kappa{=}0$ vs
gravity/$\kappa{\ne}0$ is not a physical input: it is a
theorem about the bar complex. The physical interpretation
follows from the identification of bar curvature with
perturbative genus expansions in quantum field theory.

\end{document}
